\section{Experiment 3: identifying the causal path}\label{sec:causal}

Experiment~2's headline contrast is observational. Auditing its implementation
after the fact we found two paths the design left open and the reporting did not
separate.

\textbf{The triage instruction was present during phase 2.}
\path{scripts/paper-phase2.ts:140} builds \texttt{system = SYSTEM + INVARIANT}
for the triage arm, and line 160 passes that same \texttt{system} to the phase-2
call. The randomized triage contrast (21/150 against 50/150) therefore admits a
direct route into phase-2 reasoning alongside the mediated one we reported.

\textbf{The whole phase-1 conversation was carried forward.} The message array
accumulates, so the phase-2 call saw the model's own phase-1 rationale, which in
the drifted arm frequently argued explicitly against promotional pricing.
Verified and unverified episodes differed in what the model had already
committed to in writing, not only in what evidence it held.

Experiment~3 closes both by construction. Its hypotheses, arms, and statistics
were committed to \path{workshop/PREREGISTRATION.md} before its first model call.

\subsection{A randomized-instrument replay (H9, H10)}

For each of the \ThreeAN{} Experiment-2 episodes we replay \emph{only} the
second decision as a fresh session: the triage sentence is absent from every
arm, no conversation history is carried, the memory block and its per-episode
order are reproduced from the stored episode, and a \textsc{recovered} block
holds source records for exactly the memories that episode's budget was spent
on. The phase-2 situation text is byte-identical across arms.

Triage assignment was randomized upstream and moves recovery of the corrupted
memory from \ThreeARecDrN{}/150 to \ThreeARecTrN{}/150. The replay removes the
two direct paths Experiment~2 left open: the instruction is gone from the
phase-2 session and so is the prior conversation.

It does not make the exclusion restriction automatic, and we do not claim it
does. Triage changes which records are recovered, not only whether the target's
is among them, so an instrumental estimate still assumes the other recovered
records do not themselves move the phase-2 decision. That assumption is
untestable here. We report the replay as convergent evidence and let the fully
randomized experiment in \S6.2 --- where nothing is self-selected --- carry the
causal claim.

It replicates. The corrupted-direction rate is \ThreeATrK/\ThreeATrN{}
(\ThreeATrPct\%) under triage against \ThreeADrK/\ThreeADrN{} (\ThreeADrPct\%)
without it (CMH stratified by model, $p = \ThreeACMH$). H9 is supported.

The exclusion restriction can then be checked rather than argued. Matched on
what phase~1 recovered, the arms are indistinguishable in both strata:
\ThreeARecDrK/\ThreeARecDrN{} against \ThreeARecTrK/\ThreeARecTrN{} among
episodes that recovered the caveat ($p = \ThreeARecCMH$), and
\ThreeAUnrecDrK/\ThreeAUnrecDrN{} against \ThreeAUnrecTrK/\ThreeAUnrecTrN{}
among those that did not ($p = \ThreeAUnrecCMH$). With the instruction out of
the room, arms holding the same evidence behave the same way. H10 is supported,
and reported as an equivalence observation rather than as evidence for a null.

Under that caveat the Wald estimate of the effect of recovering the corrupted
memory's source is $\IVEst$ (95\% CI $[\IVLo, \IVHi]$, stratified bootstrap
over models), within a point of the randomized estimate below. The Wald IV
estimator is a ratio of two differences and is not range-restricted, so its
interval can and here does extend past $-1$; the lower bound is an artefact of
the estimator, not a probability difference.

\paragraph{Experiment 2 was conservative, not inflated.} Removing the
carried-forward conversation \emph{raises} the corrupted-direction rate among
episodes that never recovered the caveat, from 79\% in Experiment~2 to
\ThreeAUnrecDrPctText\% here. The model's own earlier prose had been suppressing
the choice, and Experiment~2's design credited that suppression to provenance.
The confound worked against the effect we reported.

\subsection{Randomizing the provenance itself (H11)}

The replay still lets the model choose what it recovered. So we also assign it.
In \ThreeBN{} fresh single-call episodes the memory block is the drifted lineage
in both arms and two source records are carried forward: in
\textsc{carry-target} one of them is the corrupted memory's, in
\textsc{carry-other} neither is. Both arms carry exactly two records; only their
identity differs, assigned by seeded PRNG.

\begin{table}[t]
\centering\small
\caption{Study 2b. Corrupted-direction rate when the provenance carried into the
later decision is randomised rather than chosen. Both arms carry two source
records; only whether one of them is the corrupted memory's differs.}
\label{tab:carry}
\begin{tabular}{lrr}
\toprule
& corrupted memory's source carried & not carried \\
\midrule
Opus 5 & 0/25 & 14/25 \\
Sonnet 5 & 2/25 & 25/25 \\
Haiku 4.5 & 0/25 & 25/25 \\
Sol & 0/25 & 25/25 \\
Terra & 1/25 & 25/25 \\
Luna & 0/25 & 25/25 \\
\midrule
pooled & 3/150 (2\%) & 139/150 (93\%) \\
\bottomrule
\end{tabular}
\end{table}


Table~\ref{tab:carry} gives \ThreeBTgtK/\ThreeBTgtN{} against
\ThreeBOthK/\ThreeBOthN{}, a difference of \ThreeBDiff{} percentage points
(CMH $\chi^2 = \ThreeBChi$, $p \ThreeBCMH$), the same direction in all
\ThreeBAllModels{} models, with five of six choosing the corrupted direction in
every single episode where its provenance was absent. H11 is supported.

\paragraph{Randomization reaches severity, not only direction.} Experiment~2's
pre-registered harm rule returned zero because self-reported hedging saturates,
and it still does: \ThreeBNGHedged{} of the unguarded choices here omit the
uncertainty flag. But the behavioral severity measure separates cleanly under
randomization. Unguarded moves into the corrupted direction --- anything beyond
a small guarded test --- occur in \ThreeBNGOth{} episodes when the corrupted
memory's provenance is withheld and \ThreeBNGTgt{} when it is carried
(CMH $\chi^2 = \ThreeBNGChi$, $p = \ThreeBNGCMH$). Every unguarded choice in
the experiment is on the side that never saw the caveat. In the replay the same
measure gives \ThreeANonGuarded{} overall. This is the claim Experiment~2 wanted
and could not make: withholding provenance does not merely change which
direction is chosen, it changes how far the agent commits.

Two independent identification strategies --- an instrument in \S6.1, full
randomization in \S6.2 --- put the effect within a point of each other. What the
budget happened to recover is not a marker of some other difference between
episodes. It is the thing that decides.

\paragraph{Scope.} This establishes that \emph{holding} the recovered caveat
changes the later decision. It does not establish that models allocate budget in
order to secure future evidence; Experiment~1 says the opposite. The chain we
have closed is allocation $\rightarrow$ what is in context later $\rightarrow$
behavior, and contains no claim about foresight.
